\documentclass{beamer}
%
% Choose how your presentation looks.
%
% For more themes, color themes and font themes, see:
% http://deic.uab.es/~iblanes/beamer_gallery/index_by_theme.html
%
\mode<presentation>
{
  \usetheme{default}      % or try Darmstadt, Madrid, Warsaw, ...
  \usecolortheme{crane} % or try albatross, beaver, crane, ...
  \usefonttheme{structurebold}  % or try serif, structurebold, ...
  \setbeamertemplate{navigation symbols}{}
  \setbeamertemplate{caption}[numbered]
} 

\usepackage[english]{babel}
\usepackage[utf8x]{inputenc}

\usepackage{tikz-cd}

\newcommand{\C}{\mathbb{C}}
\newcommand{\R}{\mathbb{R}}
\newcommand{\cC}{\mathcal{C}}
\newcommand{\cD}{\mathcal{D}}
\renewcommand{\>}{\rangle}
\newcommand{\<}{\langle}

\title[Spectral Bounds]{Spectral Bounds for Classical \& Quantum Graph Parameters}
\author{Pawe{\l} Wocjan}
\institute{University of Central Florida}
\date{October 2019}

\begin{document}

\begin{frame}
  \titlepage
\end{frame}

% Uncomment these lines for an automatically generated outline.
\begin{frame}{Outline}
  \tableofcontents
\end{frame}

\section{Spectral graph theory}

\begin{frame}{Spectral graph theory}
\begin{itemize}
\item Spectral graph theory is a well-established branch of graph theory.

\item It studies properties of graphs relying on eigenvalues and eigenvectors of matrices associated to graphs such as their adjacency and Laplacian matrices.

\item One of the most well-known results in graph theory is the \textcolor{red}{Hoffman lower bound} on the chromatic number.
\end{itemize}
\end{frame}

\section{Graph parameters}

\begin{frame}{Graph parameters I \\ Chromatic numbers and Lov\'asz theta functions}
%\[
%\omega(G) \le \chi_v(G) = \theta'(\overline{G}) \le \chi_{sv}(G) = \theta(\overline{G}) \le \chi_{rv}(G) = \theta^+(\overline{G}) \le \chi(G).
%\]
\begin{itemize}
\item $\omega(G)$ -- clique number
\item $\chi_v(G)$ -- vector chromatic number
\item $\chi_{sv}(G)$ -- strict vector chromatic number
\item $\chi_{rv}(G)$ -- rigid vector chromatic number
\item $\chi_{\mathrm{vect}}(G)$ -- vectorial chromatic number

\medskip
\item $\theta'(\overline{G})$  -- Schrijver variant of Lov\'asz theta function
\item $\theta(\overline{G})$ -- Lov\'asz theta function 
\item $\theta^+(\overline{G})$ -- Szegedy variant of Lov\'asz theta function 

\medskip
\item These parameters, other than $\omega(G)$ can all be calculated using semi-definite programming.
\end{itemize}
\end{frame}

%

\begin{frame}{Graph parameters II}
\begin{itemize}
\item $\xi_f(G)$ -- projective rank
\item $\xi(G)$ -- orthogonal rank
\item $\xi'(G)$ -- normalized orthogonal rank
\item  $\chi_q(G)$ -- quantum chromatic number
\item $\chi_q^{(1)}(G)$ -- rank one quantum chromatic number
\item $\chi_f(G)$ -- fractional chromatic number
\item $\chi_c(G)$ -- circular chromatic number
\item $\chi(G)$ -- chromatic number
\end{itemize}
\end{frame}

%

\begin{frame}[fragile]{Hierarchies of graph parameters}

{\tiny
\[
\begin{tikzcd}
                        &                        &                                                            &                                                     & \lceil \theta^+(\overline{G}) \rceil = \chi_{\mathrm{vect}}(G)  \\
\omega(G) \ar[r] & \chi_v(G)=\theta'(\overline{G}) \ar[r] &\chi_{sv}(G) = \theta(\overline{G}) \ar[r]  & \chi_{rv}(G) = \theta^+(\overline{G}) \ar[r] \ar[ru] & \xi_f(G) 
\end{tikzcd}
\]

\medskip
\[
\begin{tikzcd}
\chi_{\mathrm{vect}}(G)   \ar[r] \ar[rd]         &  \xi(G)       \ar[to=2-3] \\ 
\xi_f(G ) \ar[to=2-2] \ar[to=1-2] \ar[to=3-2] &  \chi_q(G)  \ar[to=2-3]  & \chi_q^{(1)}(G) \ar[to=2-4] & \xi'(G) \ar[to=2-5] & \lceil \chi_c(G) \rceil = \chi(G) \\
                                                               &  \chi_f(G)     \ar[to=3-3] & \chi_c(G) \ar[to=2-5] \\
\end{tikzcd}
\]
}
Here $\longrightarrow$ is used to indicate the less or equal relationship.
\end{frame}

%

\section{Coloring}

\begin{frame}{Coloring \& chromatic number}

\begin{itemize}
\item We consider only undirected graphs without multiple edges and without loops.

\item A \textcolor{red}{$c$-coloring} of a graph $G=(V,E)$ is a collection of colors 
\[
\{ c_v : v \in V\}
\]
in $[c]=\{0,\ldots,c-1\}$ such that for all $vw\in E$
\[
c_v \neq c_w.
\]

\item The \textcolor{red}{chromatic number $\chi$} of a graph $G$ is the minimum $c$ admitting a $c$-coloring of $G$.

\item The chromatic number is NP-hard to approximate. 
\end{itemize}
\end{frame}

%

\section{Quantum coloring}

\begin{frame}{Quantum coloring, quantum chromatic number}

\begin{itemize}
\item A \textcolor{red}{quantum $c$-coloring} of a graph $G=(V,E)$ is a collection of orthogonal projectors 
\[
\{ P_{v,k} : v \in V, k \in [c] \}
\]
acting on $\C^d$ for some $d>0$ such that
\begin{itemize}
\item for all vertices $v\in V$
\[
\sum_{k\in [c]} P_{v,k} = I_d \quad \quad \mbox{completeness}
\]
\item for all edges $vw\in E$ and for all $k\in [c]$
\[
P_{v,k} P_{w,k} = 0_d \quad \quad \mbox{orthogonality}
\]
\end{itemize}

\item The \textcolor{red}{quantum chromatic number $\chi_q$} is the smallest $c$ for which the graph $G$ admits a quantum $c$-coloring.

\end{itemize}

\end{frame}

%

\begin{frame}{Classical coloring as a special case of quantum coloring}

\begin{itemize}
\item Let $\{c_v : v\in V\}$ be a collection of colors in $[c]$ that defines a classical $c$-coloring.  

\item For $v\in V$, consider the ``\emph{one-hot-encoding}''
\[
\{P_{v,k} : k \in [c]\}
\] 
of the color $c_v\in [c]$, that is,
\[
P_{v,k} = [c_ v = k],
\]
where $[\cdot ]$ denotes the Iverson bracket.

\item For instance, if $c=3$ and some vertex $v$ is colored with color $c_v=2$, then the corresponding one-hot-encoding is given by
\[
P_{v,0}=0, \quad P_{v,1}=0, \quad P_{v,2}=1.
\]
\end{itemize}

\end{frame}

% 

\begin{frame}{Classical coloring as a special case of quantum coloring}
\begin{itemize}
\item It is immediate that a classical $c$-coloring gives rise to a quantum $c$-coloring coloring in dimension $d=1$. 

\item The orthogonal projectors in the resulting quantum $c$-coloring are the scalars $0$ and $1$.
\end{itemize}
\end{frame}

%

\begin{frame}{Exponential separation between classical and quantum chromatic numbers}
\begin{itemize}
\item The \textcolor{red}{orthogonality graph}, $\Omega(n)$, has the vertex set of $\pm 1$-vectors of length $n$, with two vertices adjacent if they are orthogonal.

\item The \textcolor{blue}{quantum} chromatic number is \textcolor{blue}{linear} in $n$ (equal to $n$), whereas the \textcolor{red}{classical} chromatic number $\chi(\Omega(n))$ is \textcolor{red}{exponential} in $n$.
\end{itemize}
\end{frame}

%

\section{Hoffman bound}

\begin{frame}{Hoffman lower bound on the quantum chromatic number}

\begin{itemize}
\item Let $A=(a_{vw})\in\C^{n\times n}$ denote the adjacency matrix of a graph $G$, that is, 

\begin{itemize}
\item $a_{vw}=1$ for all $vw\in E$ 
\item $a_{vw}=0$ otherwise
\end{itemize}

\item Let $\mu_{\max}$ and $\mu_{\min}$ denote the maximum and minimum eigenvalues of $A$, respectively.

\item The \textcolor{red}{Hoffman bound} states
\[
\chi_q \ge 1 + \frac{\mu_{\max}}{|\mu_{\min|}},
\]

\end{itemize}
\end{frame}

%

\begin{frame}{Quantum coloring $\Rightarrow$ matrix equality for the adjacency matrix}

\begin{itemize}
\item We will show later that starting from a quantum $c$-coloring we can construct a unitary matrix 
\[
U\in \C^{n\times n}\otimes\C^{d\times d}
\]
such that the following \textcolor{red}{matrix equality} holds
\[
\sum_{\ell=1}^{c-1} U^\ell (-A\otimes I_d) (U^\dagger)^\ell = A \otimes I_d.
\]

\item Once we have established the above matrix equality for the adjacency matrix, we can use results from \textcolor{red}{majorization} theory to prove the Hoffman bound.
\end{itemize}

\end{frame}

%

\begin{frame}{Majorization inequality for maximum eigenvalues}

\begin{itemize}
\item Let $X$ and $Y$ be two arbitrary hermitian matrices in $\C^{m\times m}$.

\item It is well known that the maximum eigenvalues of the matrices satisfy 
\[
\mu_{\max}(X) + \mu_{\max}(Y) \ge \mu_{\max}(X + Y).
\]
\end{itemize}
\end{frame}

%

\begin{frame}{Proof of the Hoffman bound}
\begin{itemize}
\item Applying (i) the majorization bound and (ii) the invariance of the spectrum under conjugation with unitaries to the 
matrix equality for the adjacency matrix
\[
\sum_{\ell=1}^{c-1} U^\ell (-A\otimes I_d) (U^\dagger)^\ell = A \otimes I_d
\]
yields 
\[
(c-1) \mu_{\max}(-A \otimes I_d) \ge \mu_{\max}(A\otimes I_d)
\]
\item This further simplifies to
\[
(c-1) |\mu_{\min(A)}| \ge \mu_{\max}(A)
\]
which is the Hoffman lower bound.
\end{itemize}
\end{frame}

%

\section{Inertial bound}

\begin{frame}{Inertial lower bound on the quantum chromatic number}

\begin{itemize}
\item Let $(n^+, n^0, n^-)$ denote the \textcolor{red}{inertia} of a graph $G$, that is, the number of positive, zero, and negative eigenvalues of its adjacency matrix $A$.

\item The \textcolor{red}{inertial bound} states
\[
\chi_q \ge 1 + \max \left\{ \frac{n^+}{n^-}, \frac{n^-}{n^+} \right\}.
\]
\end{itemize}

\end{frame}

%

\begin{frame}{Rank inequality for positive semidefinite matrices}
\begin{itemize}
\item A hermitian matrix $X\in\C^{m\times m}$ is called \textcolor{red}{positive semidefinite} if all its eigenvalues are nonnegative.

\item Let $X,Y\in\C^{m\times m}$ be two hermitian matrices. 

\medskip
We say that $X$ is \textcolor{red}{greater or equal to} $Y$, denoted by $X\ge Y$, if their difference $X-Y$ is positive semidefinite.

\medskip
\item Assume that $X,Y\in\C^{m\times m}$ are two positive semidefinite matrices with $X\ge Y$.

\item Then, we have the rank inequality
\[
\mathrm{rank}(X) \ge \mathrm{rank}(Y).
\]
\end{itemize}
\end{frame}

%

\begin{frame}{Decomposition of the adjacency matrix in positive and negative parts}

\begin{itemize}
\item Let $\mu_1\ge \mu_2\ge \ldots \ge \mu_n$ denote the eigenvalues of the adjacency matrix $A$ and $P_1,\ldots,P_n$ the corresponding one-dimensional projectors.

\item Define 
\[
\begin{array}{lll}
A^+ = \sum_{i : \lambda_i > 0} \mu_i P_i & \quad &
P^+ = \sum_{i : \lambda_i > 0} P_i \\
\\
A^- = \sum_{i : \lambda_i < 0} |\mu_i| P_i & \quad &
P^- = \sum_{i : \lambda_i < 0} P_i 
\end{array}
\]
\item 
Then, we can write
\[
A = A^+ - A^-
\]
\item Note that
\[
n^{\pm} = \mathrm{rank}(A^{\pm})
\quad\mbox{and}\quad
A^{\pm} = \pm P^{\pm} A P^{\pm}
\]
and also $n=n^+ + n^0 + n^-$.
\end{itemize}

\end{frame}

%

\begin{frame}{Proof of the inertial bound}
\begin{itemize}
\item We start with the matrix equality
\[
\sum_{\ell=1}^{c-1} U^\ell (-A\otimes I_d) (U^\dagger)^\ell = A \otimes I_d
\]
\item Plugging in $A=A^+ - A^-$ yields
\[
\sum_{\ell=1}^{c-1} U^\ell \left((A^- - A^+\right) \otimes I_d) (U^\dagger)^\ell = (A^+ - A^-) \otimes I_d
\]
\item Multiplying both sides by $P^+$ from left and right yields
\[
\sum_{\ell=1}^{c-1} P^+ U^\ell \left((A^- - A^+\right) \otimes I_d) (U^\dagger)^\ell P^+= A^+ \otimes I_d
\]
\end{itemize}
\end{frame}

%

\begin{frame}{Proof of the inertial bound}
\begin{itemize}
\item Dropping the subtraction on the left hand side yields
\[
\sum_{\ell=1}^{c-1} P^+ U^\ell (A^-  \otimes I_d) (U^\dagger)^\ell P^+ \ge A^+ \otimes I_d
\]
\item Using the rank inequality and some elementary results yields
\[
(c-1) \cdot \mathrm{rank} (A^- \otimes I_d) \ge \mathrm{rank}(A^+ \otimes I_d)
\]
which simplifies to the inertial bound
\[
c \ge 1 + \frac{n^+}{n^-}.
\]
The other inequality $c \ge 1 + n^-/n^+$ is proved similarly.
\end{itemize}
\end{frame}

%
%
%

\begin{frame}{Proof of matrix inequality}
\begin{itemize}
\item In the following, we will prove the matrix equality
\[
\sum_{\ell=1}^{c-1} U^\ell (-A\otimes I_d) (U^\dagger)^\ell = A \otimes I_d
\]
with the help of \textcolor{red}{pinching} and \textcolor{red}{twirling}.
\end{itemize}
\end{frame}

%

\section{Pinching}

\begin{frame}{Pinching}
\begin{itemize}
\item Let $\{Q_k : k\in [c]\}$ be any collection of orthogonal projectors in $\C^{m\times m}$ that form a resolution of the identity.

\item Then, the operation $\cC$ that maps an arbitrary matrix $X\in\C^{m\times m}$ to
\[
\cC(X) = \sum_{k\in [c]} Q_k X Q_k
\]
is called \textcolor{red}{pinching}.

\item We say that the pinching $\cC$ \textcolor{red}{annihilates} $X$ if
\[
\cC(X) = 0_m.
\]
\end{itemize}
\end{frame}

%

\begin{frame}{Pinching from quantum coloring}

\begin{itemize}
\item Assume that the orthogonal projectors 
\[
P_{v,k} \quad \quad v \in V, k \in [c] 
\]
acting on $\C^d$ form a quantum $c$-coloring of $G$.

\medskip
\item Then, the block-diagonal orthogonal projectors that act on $\C^n\otimes \C^d$ and are defined by
\begin{align*}
P_k 
&= 
\bigoplus_{v\in V} P_{v,k} \\
&= 
\sum_{v\in V} |v\>\<v| \otimes P_{v,k} \quad \quad k \in [c]
\end{align*}
give rise to a pinching $\cC$.
\end{itemize}

\end{frame}

%

\begin{frame}{Pinching from quantum coloring}
\begin{itemize}
\item Due to the block diagonal structure of $P_k$ we have
\[
P_k (A\otimes I_d) P_k =
\sum_{v,w\in V} a_{vw} |v\>\<w| \otimes P_{v,k} P_{w,k}.
\]
\item Due to the orthogonality condition we have
\[
a_{vw}=1 \quad \Rightarrow \quad P_{v,k} P_{w,k} = 0_d.
\]

\item Therefore, this pinching annihilates $A\otimes I_d$, that is, 
\[
\cC(A\otimes I_d) = 0_n \otimes 0_d.
\]

\end{itemize}
\end{frame}

%

\section{Twirling}

\begin{frame}{Twirling}
\begin{itemize}
\item Let $\{U_\ell : \ell\in [c]\}$ be any collection of unitary matrices in $\C^{m\times m}$.

\item Then, the operation $\cD$ that maps an arbitrary matrix $X\in\C^{m\times m}$ to
\[
\cD(X) = \frac{1}{c} \sum_{\ell\in [c]} U_\ell X U_\ell^\dagger
\]
is called \textcolor{red}{twirling}.

\item We say that the twirling $\cD$ \textcolor{red}{annihilates} $X$ if
\[
\cD(X) = 0_m.
\]
\end{itemize}

\end{frame}

%

\begin{frame}{Twirling from pinching}

\begin{itemize}
\item Let $\{Q_k : k\in [c]\}$ be a collection of orthogonal projectors defining a pinching operation $\cC$.

\item Let $\omega = e^{2 \pi i /c}$ be a $c$th root of unity.

\item Then, the collection of unitary matrices $\{ U^\ell : \ell \in [c]\}$ consisting of the powers of the unitary $U$
\[
U = \sum_{k\in [c]} \omega^{k} Q_k,
\]
defines a twirling operation $\cD$ such that
\[
\cC(X) = \cD(X) 
\]
for all  matrices $X\in\C^{m\times m}$.

\end{itemize}

\end{frame}

%

\begin{frame}{Twirling from quantum coloring}

\begin{itemize}

\item Use the orthogonal projectors $P_{k,v}$, $v\in V, k \in [c]$, of a quantum $c$-coloring
%
to construct the corresponding block-diagonal orthogonal projectors 
\[
P_k = \sum_{v\in V} |v\>\<v| \otimes P_{v,k} \quad \quad k \in [c]
\]
acting on $\C^n\otimes \C^d$ and the unitary
\[
U = \sum_{k \in [c]} \omega^{k} P_k\,. 
\]

\item Then, the twirling map $\cD$ defined by the powers of $U$
\[
\cD(X) = \frac{1}{c} \sum_{\ell\in [c]} U^\ell X (U^\dagger)^\ell
\]
annihilates $A\otimes I_d$. This is equivalent to the matrix equality.
\end{itemize}

\end{frame}

%

\begin{frame}{Importance of twirling from quantum coloring}
\begin{itemize}
\item The construction of twirling \& pinching from quantum coloring leads to a unified way of proving the Hoffman and inertial lower bounds for $\chi_q$.

\medskip
\item Moreover, any existing or future lower bound on the minimum number of operations for pinching and twirling to annihilate $A$ is automatically a lower bound for $\chi_q$.
\end{itemize}
\end{frame}

%

\begin{frame}{Implication of the lower bounds}
\begin{itemize}
\item The spectrum of the Clebsch graph on $16$ vertices is $(5^1, 1^{10}, -3^5)$.

\item It is known that $\chi=4$.

\item The Hoffman bound implies $\chi_q \ge 3$, whereas the inertial bounds $\chi_q \ge 4$.

\item Sometimes these and other spectral lower bounds allow us to determine the quantum chromatic number. 

\item It is not known whether the quantum chromatic number is computable which is why spectral bounds are particularly useful and interesting.
\end{itemize}
\end{frame}

%

\begin{frame}{Simulating Hamiltonians}

\begin{itemize}
    \item Consider a quantum system $({\mathbb{C}^2})^{\otimes n}$ system with Hamiltonian
    \begin{equation*}
        H = \sum_{k<\ell} \sum_{a, b} c_{ab}^{k\ell} \sigma_a^k \sigma_b^\ell
    \end{equation*}
    
    \item Assume that the control operations $U\in\mathcal{U}(2)^{\otimes n}$ are available. 
    
    \item Consider of the problem of simulating the Hamiltonian $-H$ by interspersing the natural time with control operations using average Hamiltonian theory.
    
    \item The time complexity of the required control scheme depends on the chromatic number of the graph describing the coupling topology of the system Hamiltonian.
    
    \item I studied Hamiltonian simulation in my PhD and this is how I originally found my way to spectral graph theory.
\end{itemize}

\end{frame}

%

\begin{frame}{Joint work with Clive Elphick}

\begin{thebibliography}{10}

\bibitem{wocjan184}
P. Wocjan and C. Elphick, \emph{Spectral lower bounds for the orthogonal and projective ranks of a graph}, arXiv:1806.02734, Electronic Journal of Combinatorics.

\bibitem{wocjan18}
P. Wocjan and C. Elphick, \emph{Spectral lower bounds for the quantum chromatic number of a graph},  arXiv:1805.08334, Journal of Combinatorial Theory, Series A.

\bibitem{wocjan183}
P. Wocjan and C. Elphick, \emph{More Tales of Hoffman: bounds for the vector chromatic number of a graph}, arXiv:1812.02613.

\bibitem{wocjan182}
P. Wocjan and C. Elphick, \emph{An inertial upper bound for the quantum independence number of a graph},  arXiv:1808.10820.
\end{thebibliography}

\end{frame}

% 

\end{document}

%
%
%
%
%
%
%
%
%
%
%
%
%
%
%
%

\begin{frame}{Vector majorization}
\begin{itemize}
\item For an arbitrary vector $x=(x_1,\ldots,x_m)\in\R^n$, define $x^\downarrow=(x_1^\downarrow,\ldots,x_n^\downarrow)$ to be the vector whose entries are the entries of $x$ sorted in non-increasing order. 

For instance,
\[
(3.1, 1.5, 2.0)^\downarrow = (1.5, 2.0, 3.1)
\]

\item Let $x,y\in\R^m$. We say $x$ \textcolor{red}{majorizes} $y$, denoted by, $x \succeq y$, iff
\begin{align*}
\sum_{i=1}^m x_i^\downarrow &\ge \sum_{i=1}^m y_i^\downarrow \quad \mbox{for $m=1,\ldots,n-1$} \\
\sum_{i=1}^m x_i^\downarrow & \ge \sum_{i=1}^m y_i^\downarrow \quad\mbox{for $m=n$}.
\end{align*}

\end{itemize}
\end{frame}

%

\begin{frame}{Matrix majorization}
\begin{itemize}
\item Let $A\in\C^{n\times n}$ be an arbitrary hermitian matrix. Let 
\[
\mathrm{spec}(A)=(\mu_1(A),\ldots,\mu_n(A))
\]
be the spectrum of $A$, that is, the vector whose entries are the eigenvalues of $A$.

\item Let $A,B\in\C^{n\times n}$ be two arbitrary hermitian matrices. We say $A$ \textcolor{red}{majorizes} $B$, denoted by $A\succeq B$, iff the spectrum of $A$ majorizes the spectrum of $B$, that is,
\[
\mathrm{spec}(A) \succeq \mathrm{spec}(B).
\]
\end{itemize}
\end{frame}
